\documentclass[10pt,a4paper]{article}
\usepackage[margin=0.85in]{geometry}
\usepackage[T1]{fontenc}
\usepackage[utf8]{inputenc}
\usepackage{lmodern,amsmath,amssymb,graphicx,enumitem,microtype,fancyhdr,xcolor,needspace}
\usepackage[hidelinks]{hyperref}
\definecolor{lightgrey}{gray}{0.94}
\definecolor{midgrey}{gray}{0.35}
\setlength{\parindent}{0pt}
\setlength{\parskip}{5pt}
\setlist[enumerate]{itemsep=5pt,topsep=4pt}
\pagestyle{fancy}\fancyhf{}\renewcommand{\headrulewidth}{0pt}
\fancyfoot[L]{\footnotesize\color{midgrey}Arij Asad}
\fancyfoot[C]{\footnotesize\color{midgrey}Edexcel companion practice}
\fancyfoot[R]{\footnotesize\color{midgrey}\thepage}
\newcounter{question}
\newenvironment{question}[1]{\refstepcounter{question}\Needspace{5\baselineskip}\par\medskip\textbf{\thequestion.\ #1}\par}{\par\vspace{12mm}}
\begin{document}
\hypersetup{pdftitle={Statistics and Mechanics Year 1 - Original practice and solutions},pdfauthor={Arij Asad}}
\begin{minipage}[c]{0.78\textwidth}{\Large\bfseries Statistics and Mechanics Year 1}\par Original practice check\end{minipage}\hfill\begin{minipage}[c]{0.17\textwidth}\IfFileExists{PS_logo.png}{\includegraphics[width=\linewidth]{PS_logo.png}}{\rule{0pt}{16mm}}\end{minipage}

\medskip\fcolorbox{black!20}{lightgrey}{\parbox{\dimexpr\linewidth-2\fboxsep-2\fboxrule\relax}{Three short checks on probability, coded data and motion under a constant force. Include units in mechanics answers. Attempt the questions before reading the solutions. These original questions sample selected skills; they are not an official Edexcel paper.}}

\begin{question}{Binomial probability}
A biased coin has probability \(1/3\) of landing heads on each toss. Four tosses are independent. Find the probability of at least one head.
\end{question}
\begin{question}{Coding and measures of spread}
A data set has mean \(12\) and standard deviation \(3\). Every value \(x\) is replaced by \(y=2x-5\). Find the mean and standard deviation of the new data set.
\end{question}
\begin{question}{Newton's second law and constant acceleration}
A particle of mass \(3\,\mathrm{kg}\) starts from rest on a smooth horizontal surface. A constant horizontal force of \(12\,\mathrm{N}\) acts on it. Find its displacement after \(4\,\mathrm{s}\).
\end{question}
\Needspace{7\baselineskip}\textbf{Find the mistake}\par
A particle has velocity \(v=(6-2t)\,\mathrm{m\,s^{-1}}\) for \(0\leq t\leq5\), with \(t\) in seconds. A student finds a displacement of \(5\,\mathrm{m}\) and calls this the distance travelled. Find the actual distance.

\vfill\small Further practice: \href{https://maths.arijasad.com/tmua-paper-archive.html}{TMUA papers and solutions} and \href{https://maths.arijasad.com/maths-topic-practice.html}{maths topic guides}.
\newpage
{\Large\bfseries Hints and worked solutions}\par\medskip
\Needspace{8\baselineskip}\subsection*{1. Binomial probability}
\textit{Hint: The complement of at least one head is no heads.}\par
The probability of a tail on one toss is \(2/3\). Independence means that the probability of four tails is \((2/3)^4=16/81\).

Therefore \(P(\text{at least one head})=1-16/81=65/81\).

\textbf{Answer:} \(\dfrac{65}{81}\).

\Needspace{8\baselineskip}\subsection*{2. Coding and measures of spread}
\textit{Hint: A translation changes the centre but not the spread; multiplying every value scales the spread.}\par
The mean transforms in the same way as each value: \(\bar y=2\bar x-5=2(12)-5=19\).

Subtracting \(5\) does not change deviations from the mean. Multiplication by \(2\) doubles their sizes, so the standard deviation is \(2(3)=6\).

\textbf{Answer:} Mean \(19\); standard deviation \(6\).

\Needspace{8\baselineskip}\subsection*{3. Newton's second law and constant acceleration}
\textit{Hint: Find the acceleration from the horizontal resultant force, then use a constant-acceleration equation.}\par
The surface is smooth, so there is no friction. Vertically, weight and the normal reaction balance. Horizontally, \(F=ma\) gives \(12=3a\), so \(a=4\,\mathrm{m\,s^{-2}}\).

With \(u=0\) and \(t=4\), \(s=ut+\tfrac12at^2=0+\tfrac12(4)(4^2)=32\,\mathrm{m}\).

\textbf{Answer:} \(32\,\mathrm{m}\) in the direction of the force.

\Needspace{9\baselineskip}\subsection*{Find the mistake: correction}
The particle changes direction when \(6-2t=0\), at \(t=3\). Its displacement over the whole interval is indeed \([6t-t^2]_0^5=5\,\mathrm{m}\).

From \(0\) to \(3\) seconds it travels \([6t-t^2]_0^3=9\,\mathrm{m}\). From \(3\) to \(5\) seconds its displacement is \([6t-t^2]_3^5=-4\,\mathrm{m}\), so it travels a further \(4\,\mathrm{m}\).

Total distance is \(9+4=13\,\mathrm{m}\).

\Needspace{6\baselineskip}\subsection*{What to practise next}\begin{enumerate}
\item Write down the random variable and modelling assumptions before calculating a binomial probability.
\item Predict how the mean and spread should change under coding before using a calculator.
\item For mechanics, start with a force diagram, a positive direction and a list of known quantities; check units at the end.
\end{enumerate}\end{document}
